\documentclass[10pt]{article}
\usepackage[top=0.65in,bottom=0.73in,left=0.88in,right=0.88in]{geometry}
\usepackage[cmex10]{amsmath}
\usepackage[hidelinks]{hyperref}
\DeclareFontShape{OMX}{cmex}{m}{n}{<->cmex10}{}
\DeclareMathSizes{10}{10}{7}{7}
\DeclareMathSizes{9}{10}{7}{7}
\linespread{0.97}
\title{Completing a fourteen-contact anchor in dimension five}
\author{Achita Anantachina}
\date{26 November 2026}
\begin{document}
\maketitle
\begin{abstract}
We search for a 41-point kissing configuration in dimension five with
a point incident to fourteen contacts. A published anchor admits an
exact 37-point extension. We also give an integer test for mixed
rational-direction and lifted anchor points and a finite exclusion for
completion directions with coordinates bounded by two in absolute value.
\end{abstract}

\section{Construction and exact check}
A kissing configuration in dimension five is a set of unit vectors with
distinct-pair inner products at most \(1/2\). The record construction has
40 points. An exact computer-assisted upper bound of 43 is given in
Tomar's preprint. We take the fourteen integer directions from
Nyam-Osor's Table~1 and normalize them to obtain a code \(Y\subset S^3\) with
distinct-pair inner products at most \(1/3\). The pole \(e_5\) and the fourteen
points \((\sqrt{3}y/2,1/2)\), for \(y\in Y\), make up the anchor. Each
lift has inner product \(1/2\) with the pole, and two lifts have inner
product \(1/4+3\langle y_i,y_j\rangle/4\le 1/2\).

For two free integer directions \(a,b\), test \(4(a\cdot b)^2\le
\|a\|^2\|b\|^2\) whenever the dot product is positive. The pole is checked
by the same rule. For a lift from a four-dimensional integer direction
\(v\) and a five-dimensional integer direction \(a\), set \(s=\|v\|^2\),
\(n=\|a\|^2\), \(t=a_5\), and \(d=a_{1:4}\cdot v\). The constraint is
\(\sqrt{3}d/\sqrt{s}\le\sqrt n-t\). Since the right side is
nonnegative, \(d\le 0\) passes. For positive \(d\), squaring gives the
equivalent condition
\[
s(n+t^2)-3d^2 \ge 2st\sqrt n.
\]
Compare the signs of these two sides before squaring again. The check
therefore uses only integers, including on the boundary.

The archived witness gives 22 integer completion directions, and an
independent exact checker confirmed the resulting 37-point configuration.
Its purpose is to certify an extension of the given anchor. The
unrestricted lower bound remains 40.

\section{Finite search}
Enumerate all primitive, oriented integer directions \(a\) in
\([-r,r]^5\) and keep those that satisfy the fifteen anchor constraints.
Join surviving directions whenever their normalized points are
compatible. A completion is a clique. Searching backwards through
greedy independent colour classes bounds each branch by the number of
classes remaining. This is an upper bound because a clique takes at most
one vertex from each class.

At \(r=1\), the graph has 79 vertices and 2,204 edges. It has clique
number 11. At \(r=2\), it has 806 vertices and 221,136 edges. There is no
clique of size 22. These computations respect the fixed anchor. The
integer cube condition limits only the new directions.

\section{Continuous search}
We minimized a smooth maximum of the relevant dot products, then
polished the candidate using constrained optimization. A second
search also varied the fourteen horizontal directions while fixing
their fifth coordinates. All candidate coordinates were rounded to
integer directions and checked at threshold \(1/2\). The archived runs
give no 41-point counterexample. Their numerical optima give no upper
bound on what another start, anchor, or search can achieve.

\paragraph{Sources.}
\href{https://doi.org/10.1007/s00454-026-00841-x}{Cohn and Rajagopal,
Variations on five-dimensional sphere packings (2026).}
\href{https://www.mefyi.com/papers/kissing-numbers-five-seven/}{Tomar,
The kissing numbers in dimensions five, six and seven are at most 43,
76 and 132 (2026).}
\href{https://zenodo.org/records/24014108}{Nyam-Osor, Table 1 in
A three-point proof for the four-dimensional 14-point bound (2026).}
The search and manuscript used AI assistance.
\end{document}
